% !TEX root = /Users/nai1ka/Projects/thesis_template/thesis.tex
\chapter{Evaluation and Discussion}
\label{chap:eval}

This chapter evaluates the performance monitoring system presented in Chapters~\ref{chap:met} and~\ref{chap:impl}. The evaluation is organised around three questions: how much overhead the SDK adds to a host application, how accurate its measurements are compared with established profiling tools, and how reliably the regression detection pipeline identifies performance regressions. Section~\ref{sec:eval_setup} describes the validation setup. Section~\ref{sec:eval_overhead} reports the overhead introduced by the SDK. Section~\ref{sec:eval_accuracy} compares the collected metrics against reference tools. Section~\ref{sec:eval_regression} evaluates the regression detection pipeline on synthetic regressions. Section~\ref{sec:eval_discussion} interprets the results in relation to the research questions, and Section~\ref{sec:eval_limitations} discusses the limitations of the evaluation.

\section{Validation Setup}
\label{sec:eval_setup}

The evaluation was divided into three parts: a measurement of the overhead added by the SDK (Section~\ref{sec:eval_overhead}), a comparison of the collected metrics against reference tools (Section~\ref{sec:eval_accuracy}), and an evaluation of the regression detection pipeline on synthetic regressions (Section~\ref{sec:eval_regression}).

\subsection{Test Devices}

The overhead and accuracy experiments were carried out on [TODO: number] physical Android devices. The devices were selected to cover the three performance cohorts defined in Section~\ref{sec:detection}, so that the results would not be limited to a single hardware class. Table~\ref{tab:eval_devices} lists each device together with its hardware characteristics, its Android version, and the cohort assigned by the backend.

% TODO: replace the placeholder rows with the real test devices. Add or remove rows as needed.
% \begin{table}[htbp]
% \centering
% \caption{Test Devices Used in the Overhead and Accuracy Experiments}
% \label{tab:eval_devices}
% \begin{tabular}{lcccc}
% \toprule
% Device model & RAM (GB) & CPU cores & Android version & Cohort \\
% \midrule
% [TODO] & [TODO] & [TODO] & [TODO] & [TODO] \\
% [TODO] & [TODO] & [TODO] & [TODO] & [TODO] \\
% [TODO] & [TODO] & [TODO] & [TODO] & [TODO] \\
% \bottomrule
% \end{tabular}
% \end{table}

\subsection{Test Application}

The SDK was integrated into [TODO: name and a short description of the host application used for the experiments]. [TODO: describe how the application was exercised during the experiments, for example through a fixed sequence of screens and user actions, so that each run was repeatable.]

% TODO: Section 3.6 currently states that the SDK was integrated into "several open-source Android apps", and Research Question 3 refers to "real Android applications". Make sure the description above is consistent with those statements; adjust the wording in Chapter 3 and Chapter 1 if only one application or a demo application was used.

\subsection{Reference Tools}

To check the accuracy of the collected metrics, the values reported by the SDK were compared against reference tools. Android Profiler was used as the reference for CPU usage and memory usage, because it reports per-process resource consumption during a session. % TODO: add a citation for Android Profiler (for example the Android Studio profiler documentation).
System tracing provided by the Android platform was used as the reference for frame rendering time, because it records the timing of individual frames~\cite{AndroidTracing}. % TODO: confirm which tracing tool was actually used; Systrace is deprecated and has been replaced by Perfetto.
These tools are intended for use during development and are not suitable for continuous monitoring on user devices, but under controlled conditions they are accepted as reliable references.

\subsection{Synthetic Regression Injection}

The regression detection pipeline was evaluated using synthetic regressions, that is, performance degradations introduced deliberately into the host application. Three types of synthetic regression were used: an increase in CPU load, a memory leak, and blocking of the UI thread. Each type targets a different performance metric. Increased CPU load is expected to raise the CPU usage metric, a memory leak is expected to raise the memory usage metric, and UI thread blocking is expected to raise frame time and user interaction delay.

For each experiment, the application was first run in its normal state to populate the baseline window, and the synthetic regression was then enabled to populate the current window. Because the expected effect of each injected regression is known, the output of the pipeline can be compared directly against the expected result. The exact way in which each regression was injected is described in Section~\ref{sec:eval_regression}.

\section{SDK Overhead Evaluation}
\label{sec:eval_overhead}
In the first part of evaluation, the overhead that the SDK adds to a host application was measured. To do this, the same usage scenario was run in two configurations. First, the application was run with the SDK disabled, and then with the SDK enabled.
The resource usage of the application process was recorded in both configurations using the reference tools described in Section~\ref{sec:eval_setup}. To reduce the effect of random variation, the scenario was repeated [TODO: number] times on each device, and the reported values are the averages of these runs. The overhead is the difference between the two configurations.

% TODO: confirm which resources were measured. CPU and memory are included in the table below; add battery consumption and network data volume if they were measured, or remove this sentence if they were not.
The measurements covered CPU usage and memory usage. [TODO: state whether battery consumption and network data volume were also measured, and if so, describe briefly how.] Table~\ref{tab:eval_overhead} reports the results for each test device.

% TODO: replace the placeholder rows with the measured overhead values.
% \begin{table}[htbp]
% \centering
% \small
% \caption{SDK Overhead Measured on the Test Devices}
% \label{tab:eval_overhead}
% \begin{tabular}{lcccccc}
% \toprule
% \multirow{2}{*}{Device model} & \multicolumn{3}{c}{CPU usage (\%)} & \multicolumn{3}{c}{Memory usage (MB)} \\
% \cmidrule(lr){2-4} \cmidrule(lr){5-7}
%  & Without SDK & With SDK & Overhead & Without SDK & With SDK & Overhead \\
% \midrule
% [TODO] & [TODO] & [TODO] & [TODO] & [TODO] & [TODO] & [TODO] \\
% [TODO] & [TODO] & [TODO] & [TODO] & [TODO] & [TODO] & [TODO] \\
% [TODO] & [TODO] & [TODO] & [TODO] & [TODO] & [TODO] & [TODO] \\
% \bottomrule
% \end{tabular}
% \end{table}

Across the test devices, the SDK added on average [TODO: value]\% CPU usage and [TODO: value]~MB of memory. % TODO: if battery and network were measured, report the average battery overhead and the average data volume sent per device per day here.
The CPU overhead of [TODO: value]\% is [TODO: lower than / comparable to / higher than] the overhead of about 6\% reported for the Panappticon tracing system~\cite{Panappticon}. [TODO: state whether the measured overhead meets the low-overhead requirement from Section~\ref{sec:req}; this claim must be supported by the values in Table~\ref{tab:eval_overhead}.]

\section{Metric Accuracy Evaluation}
\label{sec:eval_accuracy}

The second part of the evaluation checked whether the metrics collected by the SDK are accurate. 

For each metric, the value reported by the SDK was compared against the value reported by the corresponding reference tool, measured at the same time, on the same device, and under the same usage scenario. The comparison covered CPU usage, memory usage, frame time, and startup time.  The deviation was calculated as the relative difference between the SDK value and the reference value. Table~\ref{tab:eval_accuracy} reports the results.

% TODO: replace the placeholder rows with the measured values. Report per-device values, or add a device column, if the deviation differs noticeably between devices.
% \begin{table}[htbp]
% \centering
% \caption{Accuracy of the Collected Metrics Compared with Reference Tools}
% \label{tab:eval_accuracy}
% \begin{tabular}{lccc}
% \toprule
% Metric & SDK value & Reference value & Deviation (\%) \\
% \midrule
% CPU usage & [TODO] & [TODO] & [TODO] \\
% Memory usage & [TODO] & [TODO] & [TODO] \\
% Frame time & [TODO] & [TODO] & [TODO] \\
% Startup time & [TODO] & [TODO] & [TODO] \\
% \bottomrule
% \end{tabular}
% \end{table}

The deviations were within [TODO: value]\% for all four metrics. % TODO: confirm this statement once the values are known, and revise it if any metric deviated more.
A small and consistent deviation is acceptable for this system, because the regression detection pipeline compares the relative change of a metric between the baseline and current windows rather than its absolute value (Section~\ref{sec:detection}). A constant measurement offset therefore appears in both windows and largely cancels out when the percentile shift is calculated.

\section{Regression Detection Evaluation}
\label{sec:eval_regression}

The third part of the evaluation measured how well the system detects performance regressions. The synthetic regressions described in Section~\ref{sec:eval_setup} were injected one at a time, and the output of the regression detection pipeline was compared against the expected result.

For these experiments, the baseline and current windows were shorter than the production values mentioned in Section~\ref{sec:detection} ([TODO: production window lengths, for example 7 days for the baseline window and 24 hours for the current window]). The windows were set to [TODO: experiment window lengths] so that a complete experiment could be carried out in a single session. The detection logic itself, including the P95 percentile shift, the degradation threshold, and the Mann-Whitney U test, was not changed.

A synthetic CPU regression was injected by [TODO: describe how CPU load was increased, for example by running an additional background computation on a specific screen]. The pipeline reported a P95 degradation of [TODO: value]\% for the CPU usage metric, with a Mann-Whitney U test $p$-value of [TODO: value]. [TODO: state whether the regression was confirmed.]

A synthetic memory leak was injected by [TODO: describe how the leak was introduced, for example by retaining objects so that they are never released]. The pipeline reported a P95 degradation of [TODO: value]\% for the memory usage metric, with a $p$-value of [TODO: value]. [TODO: state whether the regression was confirmed.]

A synthetic UI thread regression was injected by [TODO: describe how the UI thread was blocked, for example by performing heavy work on the main thread]. The pipeline reported a P95 degradation of [TODO: value]\% for the frame time metric, with a $p$-value of [TODO: value]. [TODO: state whether the regression was confirmed, and whether the blocking also produced ANR events.]

Table~\ref{tab:eval_regression} summarises the outcome of the three experiments.

% TODO: replace the placeholder rows with the measured values. The "Detected" column should state whether the regression was confirmed by the Mann-Whitney U test (p < 0.05).
% \begin{table}[htbp]
% \centering
% \small
% \caption{Detection Results for the Synthetic Regressions}
% \label{tab:eval_regression}
% \begin{tabular}{lcccccc}
% \toprule
% Synthetic regression & Affected metric & Baseline P95 & Current P95 & Degradation (\%) & $p$-value & Detected \\
% \midrule
% Increased CPU load & CPU usage & [TODO] & [TODO] & [TODO] & [TODO] & [TODO] \\
% Memory leak & Memory usage & [TODO] & [TODO] & [TODO] & [TODO] & [TODO] \\
% UI thread blocking & Frame time & [TODO] & [TODO] & [TODO] & [TODO] & [TODO] \\
% \bottomrule
% \end{tabular}
% \end{table}

% TODO (optional, figures intentionally omitted for now): plots of the baseline and current distributions for the detected regressions can be added here later; see figs/regression_comparison.png, figs/regression_confirmed.png, and figs/regression_frametime.png.

In addition to the three regression experiments, a control run was carried out in which the application was used normally and no regression was injected. This run measured how many false alerts the system produces under normal conditions. In the control run, the pipeline reported [TODO: number] candidate regressions and [TODO: number] confirmed regressions. [TODO: comment on whether this matches the expected result of zero confirmed regressions.] False negatives were assessed by checking whether each injected regression was confirmed: [TODO: number] of the [TODO: number] injected regressions were detected. [TODO: if any regression was missed, state which one and explain why.]

\section{Discussion of Results}
\label{sec:eval_discussion}

This section interprets the results in relation to the three research questions defined in Chapter~\ref{chap:intro}.

The first research question asks how a low-overhead, continuous monitoring system can be built so that it runs on real user devices without significantly affecting the user experience. The overhead measurements in Section~\ref{sec:eval_overhead} address this question directly. [TODO: once the overhead values are known, state whether they show that the SDK is low-overhead enough to run on real devices without a noticeable effect on the user experience.] The design choices described in Chapter~\ref{chap:impl}, in particular the two-tier buffering of metrics and the batched network transmission, were intended to keep this overhead low by reducing the number of network connections and disk writes.

The second research question asks how performance regressions can be detected accurately using statistical methods while keeping the number of false alerts low. The system addresses this with a two-stage method (Section~\ref{sec:detection}). The main filter is the rolling-window P95 percentile shift: a candidate regression is reported only when the 95th percentile of a metric increases by more than the configured threshold, so small fluctuations are ignored. The Mann-Whitney U test is then applied as a secondary check that filters out changes likely to be random noise. This secondary check is most useful for screens and cohorts with little data, where a percentile shift can occur by chance; on screens with a large amount of data, the threshold does most of the filtering. The two-stage approach differs from threshold-based tools such as Firebase Performance Monitoring, which rely on fixed limits that must be tuned manually and can produce many false alerts~\cite{FirebasePerf}. In the control run in Section~\ref{sec:eval_regression}, the pipeline produced [TODO: number] confirmed false alerts. [TODO: comment on this number as evidence that the method helps to reduce false alerts; do not claim that false alerts are eliminated.] Confirming candidate regressions with a statistical test follows the same principle as large-scale systems such as FBDetect, which validates detected regressions with hypothesis testing before alerting~\cite{FBDetect}.

The third research question asks how effective the system is at detecting regressions on real Android applications. The synthetic regression experiments show that the pipeline detected [TODO: number] of the [TODO: number] injected regressions. [TODO: interpret this result in one or two sentences.] However, synthetic regressions are simpler than the regressions that occur in real applications, so this result should be read as evidence that the detection method works as designed, not as a measurement of its effectiveness in production. The limitations of this approach are discussed in Section~\ref{sec:eval_limitations}.

\section{Limitations}
\label{sec:eval_limitations}

The evaluation has several limitations that should be considered when interpreting the results.

First, the regressions used in the evaluation were synthetic. They were injected in a controlled way and produce a clear and steady change in a single metric. Real performance regressions are often more gradual, may affect several metrics at the same time, and may appear only on specific devices or screens. The experiments therefore show that the detection method works under controlled conditions, but they do not measure how well it performs against real regressions.

Second, the system was not deployed to real users. The overhead and the detection quality were measured on a small set of test devices running a fixed usage scenario. Real usage is far more varied, and the behaviour of the system under that variety was not measured.

Third, the number of test devices and host applications was limited. [TODO: if the experiments covered only one device cohort or one application, state this explicitly here.] The results may not generalise to all device cohorts or to applications with a different structure or workload.

Fourth, the reference tools used in the accuracy evaluation are themselves approximations. Android Profiler and system tracing are accepted as references, but they have their own measurement error, and the deviation reported in Section~\ref{sec:eval_accuracy} is relative to these tools rather than to an absolute ground truth.

Fifth, the statistical test used to confirm regressions has a known weakness. On screens and cohorts with a large amount of data, almost any change becomes statistically significant, so the percentile threshold does most of the filtering and the test mainly confirms regressions that the threshold has already found. The test is most useful for screens and cohorts with little data, where a percentile shift can occur by chance. The statistical analysis in this work was kept deliberately simple, and a more detailed statistical treatment is left for future work.

Finally, the baseline and current windows used in the regression experiments were shorter than the windows intended for production use. Short windows make the experiments practical, but they do not capture long-term effects such as daily or weekly usage patterns, which could affect the baseline in a real deployment.
